\documentclass[aps,prb,reprint,amsmath,amssymb,superscriptaddress,longbibliography,floatfix]{revtex4-2}
%DIF LATEXDIFF DIFFERENCE FILE
%DIF DEL /tmp/main_prb_tdk_before_basis_order.tex   Tue May 26 22:50:05 2026
%DIF ADD main_prb_tdk.tex                           Tue May 26 22:48:53 2026

\usepackage[T1]{fontenc}
\usepackage[utf8]{inputenc}
\usepackage{graphicx}
\usepackage{bm}
\usepackage{xcolor}
\usepackage{booktabs}
\usepackage{makecell}
%DIF PREAMBLE EXTENSION ADDED BY LATEXDIFF
%DIF CTRADITIONAL PREAMBLE %DIF PREAMBLE
\RequirePackage{color}\definecolor{RED}{rgb}{1,0,0}\definecolor{BLUE}{rgb}{0,0,1} %DIF PREAMBLE
\definecolor{TDKgreen}{rgb}{0,0.45,0}
\RequirePackage[normalem]{ulem} %DIF PREAMBLE
\RequirePackage[stable]{footmisc} %DIF PREAMBLE
\DeclareOldFontCommand{\sf}{\normalfont\sffamily}{\mathsf} %DIF PREAMBLE
\providecommand{\DIFadd}[1]{{\protect\color{blue} #1}} %DIF PREAMBLE
\providecommand{\DIFdel}[1]{{\protect\color{red}#1}} %DIF PREAMBLE
%DIF SAFE PREAMBLE %DIF PREAMBLE
\providecommand{\DIFaddbegin}{} %DIF PREAMBLE
\providecommand{\DIFaddend}{} %DIF PREAMBLE
\providecommand{\DIFdelbegin}{} %DIF PREAMBLE
\providecommand{\DIFdelend}{} %DIF PREAMBLE
\providecommand{\DIFmodbegin}{} %DIF PREAMBLE
\providecommand{\DIFmodend}{} %DIF PREAMBLE
%DIF FLOATSAFE PREAMBLE %DIF PREAMBLE
\providecommand{\DIFaddFL}[1]{\DIFadd{#1}} %DIF PREAMBLE
\providecommand{\DIFdelFL}[1]{\DIFdel{#1}} %DIF PREAMBLE
\providecommand{\DIFaddbeginFL}{} %DIF PREAMBLE
\providecommand{\DIFaddendFL}{} %DIF PREAMBLE
\providecommand{\DIFdelbeginFL}{} %DIF PREAMBLE
\providecommand{\DIFdelendFL}{} %DIF PREAMBLE
%DIF COLORLISTINGS PREAMBLE %DIF PREAMBLE
\RequirePackage{listings} %DIF PREAMBLE
\RequirePackage{color} %DIF PREAMBLE
\lstdefinelanguage{DIFcode}{ %DIF PREAMBLE
%DIF DIFCODE_CTRADITIONAL %DIF PREAMBLE
  moredelim=[il][\color{red}\scriptsize]{\%DIF\ <\ }, %DIF PREAMBLE
  moredelim=[il][\color{blue}\sffamily]{\%DIF\ >\ } %DIF PREAMBLE
} %DIF PREAMBLE
\lstdefinestyle{DIFverbatimstyle}{ %DIF PREAMBLE
	language=DIFcode, %DIF PREAMBLE
	basicstyle=\ttfamily, %DIF PREAMBLE
	columns=fullflexible, %DIF PREAMBLE
	keepspaces=true %DIF PREAMBLE
} %DIF PREAMBLE
\lstnewenvironment{DIFverbatim}{\lstset{style=DIFverbatimstyle}}{} %DIF PREAMBLE
\lstnewenvironment{DIFverbatim*}{\lstset{style=DIFverbatimstyle,showspaces=true}}{} %DIF PREAMBLE
\lstset{extendedchars=\true,inputencoding=utf8}
\newcommand{\TDKin}[1]{\textcolor{blue}{#1}}
\newcommand{\TDKout}[1]{\textcolor{red}{#1}}
\newcommand{\TDKcomment}[1]{\textcolor{TDKgreen}{\textbf{[TDK:} #1\textbf{]}}}

%DIF END PREAMBLE EXTENSION ADDED BY LATEXDIFF

\begin{document}

\title{The UZH protocol: Separating and reducing Gaussian-basis and pseudopotential errors in CP2K}

\author{Hossein Mirhosseini}
\affiliation{Center for Advanced Systems Understanding (CASUS), Conrad-Schiedt-Stra{\ss}e 20, 02826 G\"orlitz, Germany}
\affiliation{Helmholtz-Zentrum Dresden-Rossendorf, Bautzner Landstra{\ss}e 400, 01328 Dresden, Germany}

\author{Tiziano M. A. M\"uller}
\affiliation{Department of Chemistry, University of Zurich, CH-8057 Zurich, Switzerland}

\author{Matthias Krack}
\affiliation{Paul Scherrer Institute, 5232 Villigen PSI, Switzerland}

\author{Thomas D. K\"uhne}
\email{tkuehne@cp2k.org}
\affiliation{Center for Advanced Systems Understanding (CASUS), Conrad-Schiedt-Stra{\ss}e 20, 02826 G\"orlitz, Germany}
\affiliation{Helmholtz-Zentrum Dresden-Rossendorf, Bautzner Landstra{\ss}e 400, 01328 Dresden, Germany}
\affiliation{Institute of Artificial Intelligence, Technische Universit\"at Dresden, Helmholtzstra{\ss}e 10, 01069 Dresden, Germany}

\author{J\"urg Hutter}
\affiliation{Department of Chemistry, University of Zurich, CH-8057 Zurich, Switzerland}

\date{\today}

\begin{abstract}
Gaussian-and-plane-wave density-functional theory combines the compactness of atom-centered orbitals with the systematic real-space representation of densities and potentials. Its accuracy, however, is controlled by two coupled approximations: the Gaussian orbital basis and the pseudopotential. We present the UZH protocol, a CP2K-based closed-loop strategy that calibrates molecularly optimized (MOLOPT) Gaussian basis sets on small molecules and validates Goedecker--Teter--Hutter (GTH) separable dual-space pseudopotentials against all-electron full-potential linearized augmented-plane-wave (FP-LAPW) equation-of-state benchmarks using the SIRIUS electronic-structure code. The protocol separates the total error into basis-set and pseudopotential components by comparing CP2K-GTH-UZH, SIRIUS-GTH-UZH, and SIRIUS FP-LAPW calculations. It extends the Materials Cloud verification strategy to a setting in which the same pseudopotential can be evaluated with a Gaussian basis and with a systematically converged plane-wave representation. This decomposition identifies whether discrepancies originate from the CP2K Gaussian basis, from the pseudopotential itself, or from their combined use. The resulting optimized UZH settings reduce the dominant errors in the original protocol and provide a practical route for improving CP2K calculations with internally generated basis sets and scalar-relativistic GTH pseudopotentials.
\end{abstract}

\maketitle

\section{Introduction}

Density-functional theory (DFT) is routinely used as a predictive tool in condensed-matter physics, chemistry, and materials science~\cite{Kohn1999RMP,Jones2015RMP}. Reproducibility studies have shown that, once exchange-correlation and physical setup are fixed, residual discrepancies between codes are mostly numerical: basis-set completeness, treatment of the core electrons, integration grids, smearing, and convergence thresholds all matter at the precision level targeted by modern verification workflows~\cite{Lejaeghere2014,Lejaeghere2016,Huber2021,Bosoni2024}. A useful protocol therefore has to do more than report agreement with a reference. It should reveal which approximation limits the result and how that approximation can be improved.

This issue is particularly important for CP2K/Quickstep~\cite{CP2K2020,CP2KMadeSimple2026}, where Kohn-Sham orbitals are represented by atom-centered Gaussian functions and densities are mapped to auxiliary plane-wave grids within the Gaussian-and-plane-wave (GPW) formalism~\cite{Lippert1997GPW}. The approach is efficient for molecules, liquids, interfaces, and extended systems, but its accuracy is governed jointly by the chosen Gaussian basis sets and norm-conserving pseudopotentials. The two are usually chosen together, and their errors are often reported as a single code-to-code deviation. Such a holistic error is useful for users, but it is insufficient for method development because it does not say whether the dominant correction should be a larger or reoptimized basis set, a revised pseudopotential, or both. This is precisely where CP2K offers a distinctive advantage: Quickstep, the SIRIUS interface, and the ATOM code are available within one coherent code ecosystem, allowing basis-set and pseudopotential errors to be analyzed and improved holistically.

Here we formulate and apply what we call the UZH protocol. The protocol combines molecularly optimized (MOLOPT) basis sets~\cite{VandeVondele2007} with the analytic Goedecker--Teter--Hutter (GTH) family of separable dual-space Gaussian pseudopotentials and its Hartwigsen--Goedecker--Hutter (HGH) relativistic extension~\cite{Goedecker1996,Hartwigsen1998,Krack2005}. Its central idea is to close a loop between molecular calibration and solid-state verification. In the first step, the Gaussian basis is optimized and calibrated against small molecules, following the MOLOPT philosophy of minimizing both energy errors and ill-conditioning of the overlap matrix~\cite{VandeVondele2007}. In the second step, the resulting CP2K settings are tested in the reproducible unary-crystal workflow of the Automated Interactive Infrastructure and Database for Computational Science (AiiDA) common workflows~\cite{Huber2021,Bosoni2024}. In the third step, the total CP2K error is decomposed by using SIRIUS both with the same GTH-UZH pseudopotentials and as a full-potential linearized augmented-plane-wave (FP-LAPW) reference~\cite{Zhang2022Computation,Zhang2023Sirius,SiriusDocs}.

The decomposition is simple but powerful. CP2K-GTH-UZH and SIRIUS-GTH-UZH differ in the orbital representation while keeping the same pseudopotential fixed; their difference isolates the Gaussian-basis component. SIRIUS-GTH-UZH and SIRIUS-FP-LAPW share the same code infrastructure but differ in the core-electron approximation; their difference isolates the pseudopotential component. The comparison of CP2K-GTH-UZH with SIRIUS-FP-LAPW gives the total practical error of the UZH protocol. Because CP2K also contains the ATOM machinery for basis-set optimization, scalar-relativistic atomic calculations, and GTH pseudopotential fitting, the same framework can diagnose and then improve the underlying numerical libraries.

\section{Theory and protocol}
\TDKcomment{The theory section and the terminology below were normalized to discuss basis-set effects before pseudopotential effects.}

\DIFdelbegin \subsection{\DIFdel{Separable dual-space Gaussian pseudopotentials}}
%DIFAUXCMD
\addtocounter{subsection}{-1}%DIFAUXCMD
%DIFDELCMD < 

%DIFDELCMD < %%%
\DIFdel{The GTH and HGH pseudopotentials are norm-conserving, angular-momentum-dependent pseudopotentials written in an analytic form that is compact in both real and reciprocal space~\cite{Goedecker1996,Hartwigsen1998}. Their local part is
}\begin{multline*}
\DIFdel{V_{\mathrm{loc}}^{\mathrm{pp}}(r) =
-\frac{Z_{\mathrm{ion}}}{r}\operatorname{erf}\left(\alpha^{\mathrm{pp}} r\right)
+ \sum_{i=1}^{4} C_i^{\mathrm{pp}}
\left(\sqrt{2}\alpha^{\mathrm{pp}} r\right)^{2i-2} }\\
\DIFdel{\times \exp\left[-\left(\alpha^{\mathrm{pp}} r\right)^2\right],
\qquad
\alpha^{\mathrm{pp}}=\frac{1}{\sqrt{2}r_{\mathrm{loc}}^{\mathrm{pp}}},
}\end{multline*}%DIFAUXCMD
\DIFdel{and the nonlocal part is written as
}\begin{displaymath}
\DIFdel{V_{\mathrm{nl}}^{\mathrm{pp}}(\mathbf r,\mathbf r')
=\sum_{lm}\sum_{ij}
\langle \mathbf r|p_i^{lm}\rangle h_{ij}^{l}
\langle p_j^{lm}|\mathbf r'\rangle .
}\end{displaymath}%DIFAUXCMD
\DIFdel{The projectors have Gaussian radial factors,
}\begin{displaymath}
\DIFdel{\langle \mathbf r|p_i^{lm}\rangle =
N_i^{l}Y^{lm}(\hat r) r^{l+2i-2}
\exp\left[-\frac{1}{2}\left(\frac{r}{r_l}\right)^2\right].
}\end{displaymath}%DIFAUXCMD
\DIFdel{This representation is well suited to GPW calculations because it can be evaluated efficiently with Gaussian orbitals and on real-space grids, while also being compatible with plane-wave calculations. The relativistic HGH generalization extends the construction across the periodic table and retains the separable analytic structure~\cite{Hartwigsen1998}. In the UZH protocol, the GTH form is optimized with CP2K's ATOM program against scalar-relativistic atomic reference calculations, including valence and semicore information as needed.
}%DIFDELCMD < 

%DIFDELCMD < %%%
\DIFdelend \subsection{MOLOPT basis sets}

MOLOPT basis sets are generally contracted Gaussian basis sets designed for density-functional calculations in gas and condensed phases~\cite{VandeVondele2007}. For atom $A$, a contracted basis function in angular-momentum channel $l$ can be written as
\begin{align}
\chi_{A l m \nu}(\mathbf r)
&=
Y_{lm}(\hat{\mathbf r}_A) r_A^{l}
\sum_{p=1}^{N_{Al}} c_{A l\nu p}
\exp\left(-\alpha_{A l p} r_A^2\right),
\label{eq:contracted_basis}\\
r_A&=|\mathbf r-\mathbf R_A| ,
\end{align}
where $\nu$ labels contractions that share a compact primitive exponent set $\{\alpha_{A l p}\}$ and differ through the coefficients $c_{A l\nu p}$. Valence and polarization flexibility are therefore introduced mainly through the contraction pattern, while the same primitive pool can be kept sufficiently compact for condensed-phase calculations.

The numerical stability of a candidate basis is measured with the molecular overlap matrix,
\begin{equation}
S_{\mu\nu}^{(M)}=\langle \chi_\mu|\chi_\nu\rangle_M,
\qquad
\kappa_B(M)=
\frac{\lambda_{\max}\!\left[S_B^{(M)}\right]}{\lambda_{\min}\!\left[S_B^{(M)}\right]} ,
\label{eq:overlap_condition}
\end{equation}
where $\kappa_B(M)$ is the condition number for basis level $B$ and molecule $M$. The MOLOPT optimization target balances accuracy and numerical stability,
\begin{equation}
\Omega =
\sum_{B\in\mathcal B}\sum_{M\in\mathcal T}
\left[
\sum_{q\in\mathcal Q} w_q
\left(
\frac{q_{B,M}-q_M^{\mathrm{ref}}}{s_q}
\right)^2
+ \gamma\ln\kappa_B(M)
\right],
\label{eq:molopt}
\end{equation}
where $\mathcal B$ is the set of related basis levels, $\mathcal T$ is the molecular training set, and $\mathcal Q$ contains the energy and property targets used during calibration. The scale factors $s_q$ make heterogeneous quantities comparable, and the weights $w_q$ control their relative importance. The logarithmic condition-number term is essential: diffuse primitives improve weak interactions and response properties, but unconstrained diffuse functions can create near-linear dependencies. In the MOLOPT construction, diffuse primitives are included within contractions, providing flexibility without making the overlap matrix numerically fragile.

The UZH basis optimization follows this logic but makes the calibration explicit. Small molecules are used first because they probe chemically diverse bonding situations at modest cost, permit comparison to all-electron Gaussian references, and expose failures in geometries, dipoles, polarizabilities, and vibrational frequencies before the settings are transferred to solids. This molecular stage is not the final validation target; rather, it establishes a chemically balanced basis library that can then be assessed in a reproducible periodic workflow.

\DIFaddbegin \subsection{\DIFadd{Separable dual-space Gaussian pseudopotentials}}

\DIFadd{The GTH and HGH pseudopotentials are norm-conserving, angular-momentum-dependent pseudopotentials written in an analytic form that is compact in both real and reciprocal space~\cite{Goedecker1996,Hartwigsen1998}. Their local part is
}\begin{multline}
\DIFadd{V_{\mathrm{loc}}^{\mathrm{pp}}(r) =
-\frac{Z_{\mathrm{ion}}}{r}\operatorname{erf}\left(\alpha^{\mathrm{pp}} r\right)
+ \sum_{i=1}^{4} C_i^{\mathrm{pp}}
\left(\sqrt{2}\alpha^{\mathrm{pp}} r\right)^{2i-2} }\\
\DIFadd{\times \exp\left[-\left(\alpha^{\mathrm{pp}} r\right)^2\right],
\qquad
\alpha^{\mathrm{pp}}=\frac{1}{\sqrt{2}r_{\mathrm{loc}}^{\mathrm{pp}}},
}\end{multline}
\DIFadd{and the nonlocal part is written as
}\begin{equation}
\DIFadd{V_{\mathrm{nl}}^{\mathrm{pp}}(\mathbf r,\mathbf r')
=\sum_{lm}\sum_{ij}
\langle \mathbf r|p_i^{lm}\rangle h_{ij}^{l}
\langle p_j^{lm}|\mathbf r'\rangle .
}\end{equation}
\DIFadd{The projectors have Gaussian radial factors,
}\begin{equation}
\DIFadd{\langle \mathbf r|p_i^{lm}\rangle =
N_i^{l}Y^{lm}(\hat r) r^{l+2i-2}
\exp\left[-\frac{1}{2}\left(\frac{r}{r_l}\right)^2\right].
}\end{equation}
\DIFadd{This representation is well suited to GPW calculations because it can be evaluated efficiently with Gaussian orbitals and on real-space grids, while also being compatible with plane-wave calculations. The relativistic HGH generalization extends the construction across the periodic table and retains the separable analytic structure~\cite{Hartwigsen1998}. In the UZH protocol, the GTH form is optimized with CP2K's ATOM program against scalar-relativistic atomic reference calculations, including valence and semicore information as needed.
}

\DIFaddend \subsection{Error decomposition}

Let $E_{\mathrm{C}}^{\mathrm{GTH}}(V)$ denote CP2K-GTH-UZH energies, $E_{\mathrm{S}}^{\mathrm{GTH}}(V)$ SIRIUS-GTH-UZH energies, and $E_{\mathrm{S}}^{\mathrm{AE}}(V)$ SIRIUS-FP-LAPW all-electron (AE) energies for the same structure and volume. The total deviation of the practical CP2K protocol is
\begin{equation}
\Delta_{\mathrm{tot}}(V)=
E_{\mathrm{C}}^{\mathrm{GTH}}(V)-E_{\mathrm{S}}^{\mathrm{AE}}(V).
\end{equation}
It can be decomposed as
\begin{align}
\Delta_{\mathrm{tot}}(V)
&=\Delta_{\mathrm{basis}}(V)+\Delta_{\mathrm{pp}}(V),\label{eq:decomposition}\\
\Delta_{\mathrm{basis}}(V)
&=E_{\mathrm{C}}^{\mathrm{GTH}}(V)-E_{\mathrm{S}}^{\mathrm{GTH}}(V),\\
\Delta_{\mathrm{pp}}(V)
&=E_{\mathrm{S}}^{\mathrm{GTH}}(V)-E_{\mathrm{S}}^{\mathrm{AE}}(V).
\end{align}
The first term measures the CP2K Gaussian-basis error for a fixed pseudopotential; the second term measures the pseudopotential error in a systematic basis. In practice, the comparisons are performed not at a single volume but over equation-of-state curves. We use the $\varepsilon$ metric of the verification project~\cite{Bosoni2024},
\begin{equation}
\varepsilon(a,b)=
\sqrt{
\frac{\sum_i \left[E_a(V_i)-E_b(V_i)\right]^2}{\sqrt{
\sum_i\left[E_a(V_i)-\langle E_a\rangle\right]^2
\sum_i\left[E_b(V_i)-\langle E_b\rangle\right]^2}}},
\label{eq:epsilon}
\end{equation}
where the volumes $V_i$ are sampled around the equilibrium region. This dimensionless metric is sensitive to equation-of-state differences while remaining comparable across chemically distinct elements.

\begin{figure}
\includegraphics[width=\linewidth]{comparison_b.png}
\caption{\label{fig:decomposition}
Error decomposition used in the UZH protocol. CP2K-GTH-UZH compared with SIRIUS-FP-LAPW gives the total practical error. CP2K-GTH-UZH compared with SIRIUS-GTH-UZH isolates the Gaussian-basis contribution. SIRIUS-GTH-UZH compared with SIRIUS-FP-LAPW isolates the pseudopotential contribution.}
\end{figure}

\section{Computational methods}

\subsection{Molecular calibration}

The molecular calibration follows the MOLOPT strategy of optimizing transferable, numerically well-conditioned Gaussian basis sets for chemical accuracy across molecules and condensed phases~\cite{VandeVondele2007,MullerThesis2024}. We benchmarked double-zeta valence plus polarization (DZVP), triple-zeta valence plus polarization (TZVP), and triple-zeta valence plus two polarization shells (TZV2P) UZH-MOLOPT basis sets on a database of small molecules containing diverse bonding motifs. Reference calculations were performed with Gaussian~16~\cite{Gaussian16} using all-electron def2-QZVP-quality settings, tight self-consistent-field convergence, and dense numerical integration grids. The CP2K calculations used matching all-electron and GTH-UZH pseudopotential setups to separate basis and pseudopotential effects at the molecular level. Geometries, bond angles, dihedral angles, dipole moments, polarizabilities, and vibrational frequencies were compared. Additional details and property-specific statistics are given in the Supplemental Material~\cite{SuppMat}.

\subsection{Periodic verification workflow}

Periodic benchmarks were generated with the unary-crystal verification workflow of AiiDA common workflows~\cite{Huber2021,Bosoni2024}. For each element, face-centered cubic, body-centered cubic, simple cubic, and diamond prototypes were considered where applicable. Seven volumes were sampled from 94\% to 106\% around the reference volume and fitted to a Birch-Murnaghan equation of state. The same structural inputs and workflow logic were used for CP2K-GTH-UZH, SIRIUS-GTH-UZH, and SIRIUS-FP-LAPW.

CP2K/Quickstep calculations used the GPW method with UZH-MOLOPT basis sets and the corresponding GTH-UZH pseudopotentials. SIRIUS-GTH-UZH calculations used the same GTH-UZH pseudopotentials in a plane-wave representation, thereby removing the Gaussian-basis approximation while retaining the pseudopotential. SIRIUS-FP-LAPW calculations served as the all-electron reference. SIRIUS is a domain-specific electronic-structure library that implements pseudopotential plane-wave and full-potential linearized augmented-plane-wave methods and is designed for parallel and accelerator-enabled calculations~\cite{Zhang2022Computation,Zhang2023Sirius,SiriusDocs}. In the FP-LAPW mode, the unit cell is partitioned into atom-centered augmentation spheres and an interstitial region; core and valence electrons are treated explicitly, and the density and effective potential are not constrained to a shape approximation. This makes SIRIUS-FP-LAPW an appropriate reference for identifying the residual core approximation in the pseudopotential calculations.

The relativistic treatment was kept consistent with the scalar-relativistic GTH/HGH construction. In the all-electron SIRIUS reference, the full-potential species definitions specify the relativistic treatment of core and valence states; in the pseudopotential calculations, the corresponding scalar-relativistic information is folded into the GTH-UZH potential. Spin-orbit effects are therefore not part of the error decomposition considered here. This choice is deliberate: the objective is to separate Gaussian-basis incompleteness from scalar-relativistic pseudopotential transferability for the CP2K settings used in production GPW calculations. The availability of SIRIUS inside CP2K makes this decomposition especially useful because CP2K can be used both as the production GPW code and as an interface to a systematic plane-wave or all-electron reference.

\subsection{CP2K ATOM and closed-loop optimization}

The optimization loop uses CP2K's ATOM code. Inspection of the CP2K source and input reference shows that the ATOM input layer supports energy calculations, basis optimization, and pseudopotential optimization as distinct run types; scalar-relativistic atomic reference calculations can be generated with zeroth-order regular approximation (ZORA)-type and Douglas--Kroll--Hess Hamiltonians; and GTH pseudopotentials can be read, printed, fitted, and exported~\cite{CP2K2020,CP2KMadeSimple2026}. The \DIFaddbegin \DIFadd{basis optimizer uses template, work, and output basis files and allows the condition number of the overlap matrix to enter the training objective. The }\DIFaddend pseudopotential optimizer constructs all-electron and pseudopotential atomic references with explicit core, valence, and semicore occupations, evaluates orbital energies and radial charges, and uses a Powell-type derivative-free optimization to fit the compact GTH parameter set~\cite{Powell1964}. \DIFdelbegin \DIFdel{The basis optimizer uses template, work, and output basis files and allows the condition number of the overlap matrix to enter the training objective. }\DIFdelend These implementation details are important because they make CP2K not merely the code under test, but also the code that can generate revised scalar-relativistic basis and pseudopotential libraries after the error decomposition has identified the limiting approximation.

\section{Results and discussion}

\subsection{Molecular basis calibration}

The molecular benchmark confirms the expected systematic behavior of the UZH-MOLOPT sequence. Bond-length errors decrease from DZVP to TZVP and TZV2P, with the median relative error remaining small across the database. Carbon-containing bonds are particularly stable, while O- and F-containing compounds are more sensitive to polarization and diffuse flexibility in the smaller basis sets. The TZV2P level removes most of these systematic trends.

\begin{figure}
\includegraphics[width=\linewidth]{TM_mol_1.png}
\caption{\label{fig:bonds}
Distribution of relative bond-length errors for the molecular calibration set. Outliers are omitted from the box plot to emphasize the median and interquartile behavior. The systematic narrowing from DZVP to TZV2P indicates controlled basis-set convergence for local structural properties.}
\end{figure}

Angular properties follow the same pattern. Nonlinear bond angles and dihedral angles are more sensitive than bond lengths to subtle changes in the electronic structure, but the hierarchy of basis quality remains clear. Dipole moments and polarizabilities provide a stronger test because they probe the response of the density rather than only the optimized geometry. The larger variance observed for polarizabilities is consistent with the need for diffuse flexibility, but the absence of catastrophic outliers at the TZV2P level indicates that the condition-number constraint in Eq.~(\ref{eq:molopt}) is effective.

\begin{figure*}
\includegraphics[width=0.49\textwidth]{TM_mol_2.png}
\includegraphics[width=0.49\textwidth]{TM_mol_3.png}
\caption{\label{fig:molecular_response}
Molecular calibration beyond bond lengths. Left: relative errors for nonlinear and dihedral angles. Right: relative errors for dipole moments and average molecular polarizabilities. Response properties are more demanding than equilibrium bond lengths, but they retain the systematic basis-set improvement needed before transferring the basis to periodic verification.}
\end{figure*}

\subsection{Total error of the original UZH settings}

Figure~\ref{fig:total_old} shows the total practical error of the original CP2K-GTH-UZH settings relative to the SIRIUS-FP-LAPW reference. Several elements exhibit $\varepsilon>1$ for at least two unary structures, including noble gases, selected transition metals, and heavier main-group elements. This map is the error a user would observe if CP2K were treated as a black box. It is therefore the most direct measure of the protocol's practical accuracy, but by itself it does not reveal the origin of the discrepancy.

\begin{figure*}
\includegraphics[width=\textwidth]{periodic-table-unaries-CP2K-UZH-vs-SIRIUS-LAPW-epsilon.png}
\caption{\label{fig:total_old}
Total equation-of-state error of the original CP2K-GTH-UZH settings relative to SIRIUS-FP-LAPW, expressed with the $\varepsilon$ metric of Eq.~(\ref{eq:epsilon}). This comparison contains both Gaussian-basis and pseudopotential errors.}
\end{figure*}

\begin{table*}
\caption{\label{tab:classification}
Qualitative classification of dominant error sources in the original UZH settings. The classification is based on the pairwise comparisons in Figs.~\ref{fig:pp_basis_old} and \ref{fig:total_old}.}
\begin{ruledtabular}
\begin{tabular}{lll}
Dominant source & Diagnostic signature & Representative elements \\
\hline
\DIFdelFL{Pseudopotential} & \DIFdelFL{Large SIRIUS-GTH-UZH vs SIRIUS-FP-LAPW error} & \DIFdelFL{Ca, Cr, Pd, Rh, Ru, Tc} \\
Basis set & Large CP2K-GTH-UZH vs SIRIUS-GTH-UZH error & He, Ne, Ar, Kr, Xe, Rn, Ba \\
\DIFaddbeginFL \DIFaddFL{Pseudopotential }& \DIFaddFL{Large SIRIUS-GTH-UZH vs SIRIUS-FP-LAPW error }& \DIFaddFL{Ca, Cr, Pd, Rh, Ru, Tc }\\
\DIFaddendFL Mixed & Both components visible & Hg and selected heavy elements \\
\end{tabular}
\end{ruledtabular}
\end{table*}

\subsection{Separating basis and pseudopotential errors}

The decomposition in Eq.~(\ref{eq:decomposition}) reveals three distinct regimes. First, \DIFdelbegin \DIFdel{for elements such as Ca, Cr, Pd, Rh, Ru, and Tc, SIRIUS-GTH-UZH differs substantially from SIRIUS-FP-LAPW, while CP2K-GTH-UZH remains close to SIRIUS-GTH-UZH. These cases are pseudopotential limited. Improving the Gaussian basis alone would not remove the dominant error because the same discrepancy is already present in the plane-wave representation of the pseudopotential.
}%DIFDELCMD < 

%DIFDELCMD < %%%
\DIFdel{Second, the }\DIFdelend \DIFaddbegin \DIFadd{the }\DIFaddend noble gases and selected elements such as Ba are basis limited. SIRIUS-GTH-UZH agrees well with SIRIUS-FP-LAPW, indicating that the GTH-UZH pseudopotential is transferable for these tests, but CP2K-GTH-UZH deviates from SIRIUS-GTH-UZH. These cases point to missing diffuse or polarization flexibility in the Gaussian basis for weakly bound or highly compressible solids. This observation is physically plausible because noble-gas equation-of-state curves are shallow and therefore sensitive to small basis-set and basis-set-superposition effects.

\DIFaddbegin \DIFadd{Second, for elements such as Ca, Cr, Pd, Rh, Ru, and Tc, SIRIUS-GTH-UZH differs substantially from SIRIUS-FP-LAPW, while CP2K-GTH-UZH remains close to SIRIUS-GTH-UZH. These cases are pseudopotential limited. Improving the Gaussian basis alone would not remove the dominant error because the same discrepancy is already present in the plane-wave representation of the pseudopotential.
}

\DIFaddend Third, some heavy elements show mixed behavior. In these cases both pairwise comparisons are non-negligible, and the improvement strategy must combine pseudopotential refitting with basis-set extension or reoptimization.

\begin{figure*}
\includegraphics[width=0.49\textwidth]{periodic-table-unaries-SIRIUS-UZH-vs-SIRIUS-LAPW-epsilon.png}
\includegraphics[width=0.49\textwidth]{periodic-table-unaries-CP2K-UZH-vs-SIRIUS-UZH-epsilon.png}
\caption{\label{fig:pp_basis_old}
Error decomposition for the original UZH settings. \DIFdelbeginFL \DIFdelFL{Left}\DIFdelendFL \DIFaddbeginFL \DIFaddFL{Right}\DIFaddendFL : \DIFdelbeginFL \DIFdelFL{SIRIUS-GTH-UZH }\DIFdelendFL \DIFaddbeginFL \DIFaddFL{CP2K-GTH-UZH }\DIFaddendFL relative to \DIFdelbeginFL \DIFdelFL{SIRIUS-FP-LAPW}\DIFdelendFL \DIFaddbeginFL \DIFaddFL{SIRIUS-GTH-UZH}\DIFaddendFL , which isolates the \DIFdelbeginFL \DIFdelFL{pseudopotential }\DIFdelendFL \DIFaddbeginFL \DIFaddFL{Gaussian-basis }\DIFaddendFL error. \DIFdelbeginFL \DIFdelFL{Right}\DIFdelendFL \DIFaddbeginFL \DIFaddFL{Left}\DIFaddendFL : \DIFdelbeginFL \DIFdelFL{CP2K-GTH-UZH }\DIFdelendFL \DIFaddbeginFL \DIFaddFL{SIRIUS-GTH-UZH }\DIFaddendFL relative to \DIFdelbeginFL \DIFdelFL{SIRIUS-GTH-UZH}\DIFdelendFL \DIFaddbeginFL \DIFaddFL{SIRIUS-FP-LAPW}\DIFaddendFL , which isolates the \DIFdelbeginFL \DIFdelFL{Gaussian-basis }\DIFdelendFL \DIFaddbeginFL \DIFaddFL{pseudopotential }\DIFaddendFL error.}
\end{figure*}

\subsection{Improved UZH settings}

The classification in Table~\ref{tab:classification} was used to revise the settings in a targeted way. \DIFaddbegin \DIFadd{Basis-limited elements were treated by extending or reoptimizing the UZH-MOLOPT basis while retaining the condition-number penalty that prevents unstable diffuse functions. }\DIFaddend Pseudopotential-limited elements were treated by refitting the GTH parameters with the CP2K ATOM code, including scalar-relativistic atomic reference information and semicore states where appropriate. \DIFdelbegin \DIFdel{Basis-limited elements were treated by extending or reoptimizing the UZH-MOLOPT basis while retaining the condition-number penalty that prevents unstable diffuse functions. }\DIFdelend This targeted route differs from a brute-force convergence exercise: it changes only the component that the error decomposition identifies as limiting.

Figure~\ref{fig:improved} shows the resulting total error of CP2K-GTH-UZH-new relative to SIRIUS-FP-LAPW and the residual decomposition. The improved settings reduce the dominant outliers of the original protocol and make the remaining discrepancies more interpretable. In particular, noble-gas and Ba errors are reduced by basis-set changes, while the transition-metal cases respond to pseudopotential refitting. The resulting UZH protocol is therefore not a fixed table of parameters but a reproducible recipe for diagnosing, optimizing, and validating CP2K settings.

\begin{figure*}
\includegraphics[width=0.49\textwidth]{periodic-table-unaries-CP2K-UZH-new-vs-SIRIUS-LAPW-epsilon.png}
\includegraphics[width=0.49\textwidth]{periodic-table-unaries-SIRIUS-UZH-new-vs-SIRIUS-LAPW-epsilon.png}\\
\includegraphics[width=0.49\textwidth]{periodic-table-unaries-CP2K-UZH-new-vs-SIRIUS-UZH-epsilon.png}
\includegraphics[width=0.49\textwidth]{periodic-table-unaries-SIRIUS-UZH-new-vs-SIRIUS-UZH-epsilon.png}
\caption{\label{fig:improved}
Verification of the improved UZH settings. Top left: total CP2K-GTH-UZH-new error relative to SIRIUS-FP-LAPW. \DIFaddbeginFL \DIFaddFL{Bottom left: residual CP2K Gaussian-basis error. }\DIFaddendFL Top right: residual pseudopotential error after improvement. Bottom \DIFdelbeginFL \DIFdelFL{left: residual CP2K Gaussian-basis error. Bottom }\DIFdelendFL right: direct comparison of new and original SIRIUS-GTH-UZH pseudopotential settings.}
\end{figure*}

\section{Implications for CP2K verification}

The UZH protocol turns CP2K verification into a diagnostic and constructive workflow. The comparison with SIRIUS-FP-LAPW gives the total accuracy relevant for production calculations. The intermediate SIRIUS-GTH-UZH calculation tells whether the Gaussian basis or the pseudopotential is responsible. The CP2K ATOM code then supplies the internal machinery for improvement: \DIFaddbegin \DIFadd{basis optimization with condition-number control, }\DIFaddend scalar-relativistic atomic references, \DIFaddbegin \DIFadd{and }\DIFaddend GTH parameter fitting\DIFdelbegin \DIFdel{, and basis optimization with condition-number control}\DIFdelend . This is particularly valuable for users who need reliable calculations across molecules and solids, because the protocol connects molecular transferability to periodic equation-of-state precision.

The protocol also clarifies what should be considered converged. Increasing the GPW grid cutoff alone cannot remove a Gaussian-basis error. Likewise, replacing a basis set cannot repair a pseudopotential whose scattering or semicore transferability is insufficient. Conversely, once the SIRIUS-GTH-UZH and SIRIUS-FP-LAPW curves agree, the remaining CP2K discrepancy is a basis issue that can be addressed without changing the pseudopotential. This separation is the practical advantage of combining CP2K and SIRIUS in a single AiiDA-based verification framework.

\section{Conclusions}

We have presented the UZH protocol for constructing, diagnosing, and improving CP2K basis-set and pseudopotential settings. The protocol combines molecularly optimized MOLOPT basis sets with GTH separable dual-space Gaussian pseudopotentials, calibrates the basis on small molecules, and validates the resulting settings through AiiDA unary-crystal equation-of-state workflows. By comparing CP2K-GTH-UZH, SIRIUS-GTH-UZH, and SIRIUS-FP-LAPW, the total CP2K error can be decomposed into Gaussian-basis and pseudopotential components. The same CP2K infrastructure used for production calculations also provides the ATOM and basis-optimization tools required to improve the limiting component. The resulting workflow provides a reproducible path toward systematically improvable CP2K simulations across molecules and condensed phases.

\begin{acknowledgments}
The authors gratefully acknowledge computing time on the high-performance computer Noctua2 at the National High-Performance Computing (NHR) Center Paderborn Center for Parallel Computing (PC2). This center is jointly supported by the Federal Ministry of Research, Technology and Space and the state governments participating in the NHR joint funding program. Additional support from the CP2K, SIRIUS, and AiiDA developer communities is gratefully acknowledged.
\end{acknowledgments}

\section*{Data Availability}

The curated release repository for the UZH protocol is available at \url{https://github.com/DCM-Uni-Paderborn/UZH-protocol}. It contains the optimized CP2K basis and pseudopotential files, manuscript figures, provenance manifests, and the release structure for the CP2K-GTH-UZH, SIRIUS-GTH-UZH, and SIRIUS-FP-LAPW workflow data. Complete workflow inputs, parsed outputs, and post-processing scripts will be deposited in the same repository before final publication.

\section*{Author Contributions}

Hossein Mirhosseini: conceptualization, data curation, formal analysis, methodology, software, visualization, writing--original draft, writing--review and editing. Tiziano M. A. M\"uller: data curation, methodology, software, validation, writing--review and editing. Matthias Krack: methodology, basis-set and pseudopotential curation, validation, writing--review and editing. Thomas D. K\"uhne: conceptualization, methodology, supervision, funding acquisition, writing--original draft, writing--review and editing. J\"urg Hutter: conceptualization, methodology, supervision, writing--review and editing.

\bibliography{bib_prb_tdk}

\end{document}
